\documentclass[10pt,aspectratio=169]{beamer}
\usetheme{Madrid}
\usecolortheme{default}

% Tokyo Night Inspired Palette for Beamer
\definecolor{tokyoDark}{HTML}{1a1b26}
\definecolor{tokyoBlue}{HTML}{7aa2f7}
\definecolor{tokyoPurple}{HTML}{bb9af7}
\definecolor{tokyoTeal}{HTML}{7dcfff}
\definecolor{tokyoText}{HTML}{3b4261}

\setbeamercolor{structure}{fg=tokyoBlue}
\setbeamercolor{palette primary}{bg=tokyoDark,fg=white}
\setbeamercolor{palette secondary}{bg=tokyoBlue,fg=tokyoDark}
\setbeamercolor{palette tertiary}{bg=tokyoPurple,fg=white}
\setbeamercolor{title}{bg=tokyoDark,fg=tokyoBlue}

\usepackage[utf8]{inputenc}
\usepackage[portuguese]{babel}
\usepackage{graphicx}
\usepackage{booktabs}
\usepackage{amsmath}

\title[Seminário IntraSOM]{Mapas de Kohonen com IntraSOM}
\subtitle{Synthetic Control, Visualização, Clusters e Extensão Textual}
\author[Cauã Vitor]{Cauã Vitor}
\institute[UFRN]{Universidade Federal do Rio Grande do Norte (UFRN)}
\date{Julho de 2026}

\begin{document}

\begin{frame}
  \titlepage
\end{frame}

\begin{frame}{1. Introdução e Objetivo}
  \begin{itemize}
    \item \textbf{Objetivo:} Explorar de forma não-supervisionada a capacidade dos Mapas de Kohonen (SOM) de projetar dados de alta dimensionalidade em grades bidimensionais hexagonais.
    \item \textbf{Análise Visual:} Avaliar a preservação da topologia e das relações de vizinhança na malha 2D.
    \item \textbf{Frentes de Teste:}
      \begin{enumerate}
        \item \textbf{Synthetic Control Chart (UCI):} 600 séries temporais, 60 atributos, 6 classes de tendência.
        \item \textbf{Notícias (20 Newsgroups):} 400 documentos, 4 categorias (Graphics, Baseball, Mideast, Space).
      \end{enumerate}
    \item \textbf{Ferramenta Principal:} Biblioteca \texttt{IntraSOM/USP} integrada a um Dashboard Web Interativo.
  \end{itemize}
\end{frame}

\begin{frame}{2. Motivação Científica}
  \begin{itemize}
    \item \textbf{Desafio:} Como visualizar, clusterizar e analisar dados complexos sem depender de rótulos de treinamento?
    \item \textbf{Limitação de Métodos Tradicionais:}
      \begin{itemize}
        \item PCA: Projeção linear, perde relacionamentos não-lineares complexos.
        \item t-SNE / UMAP: Excelente separação, mas não cria protótipos de dados (neurônios) representativos nem mantém distâncias de forma paramétrica contínua.
      \end{itemize}
    \item \textbf{Vantagem do SOM:} Combina redução de dimensionalidade e clusterização, gerando uma grade ordenada de representações (codebooks) que facilita a interpretabilidade humana.
  \end{itemize}
\end{frame}

\begin{frame}{3. O Conceito do SOM (Self-Organizing Maps)}
  \begin{itemize}
    \item \textbf{Arquitetura:} Grade bidimensional de neurônios, cada um associado a um vetor de pesos (codebook) de dimensão igual à dos dados de entrada.
    \item \textbf{Treinamento Não-Supervisionado:}
      \begin{enumerate}
        \item Apresentação de uma amostra de entrada $\vec{x}$.
        \item Seleção do neurônio vencedor (BMU - Best Matching Unit) via distância Euclidiana mínima:
          \[ d(\vec{x}, \vec{w}_i) = \min_j d(\vec{x}, \vec{w}_j) \]
        \item Atualização dos pesos do BMU e de seus vizinhos na grade hexagonal:
          \[ \vec{w}_j(t+1) = \vec{w}_j(t) + \eta(t) h_{ij}(t) [\vec{x} - \vec{w}_j(t)] \]
      \end{enumerate}
    \item \textbf{Topologia Hexagonal:} Oferece distância equidistante para todos os 6 vizinhos imediatos na grade, melhor que a grade retangular.
  \end{itemize}
\end{frame}

\begin{frame}{4. A Biblioteca IntraSOM (USP)}
  \begin{itemize}
    \item \textbf{Origem:} Biblioteca Python desenvolvida pelo grupo InTRA da USP para treinamento, visualização e avaliação quantitativa de Mapas SOM.
    \item \textbf{Recursos Chave:}
      \begin{itemize}
        \item Geração nativa da U-Matrix (Matriz de Distâncias Unificada).
        \item Integração com algoritmos de agrupamento de múltiplos níveis.
        \item Tratamento robusto para valores nulos (missing data) e imputação baseada no SOM.
        \item Cálculo automatizado de erros topológicos e de quantização.
      \end{itemize}
  \end{itemize}
\end{frame}

\begin{frame}{5. Dataset Principal: Synthetic Control Chart (UCI)}
  \begin{itemize}
    \item \textbf{Especificações:} 600 séries temporais sintéticas de 60 atributos cada (comportamento temporal).
    \item \textbf{Estrutura de Classes:}
      \begin{table}
        \centering
        \begin{tabular}{lcl}
          \toprule
          Classe & Intervalo & Comportamento Típico \\
          \midrule
          Normal & 1 -- 100 & Estável, flutuações pequenas ao redor da média \\
          Cyclic & 101 -- 200 & Oscilações senoidais regulares \\
          Increasing Trend & 201 -- 300 & Tendência contínua de subida \\
          Decreasing Trend & 301 -- 400 & Tendência contínua de descida \\
          Upward Shift & 401 -- 500 & Salto abrupto para cima \\
          Downward Shift & 501 -- 600 & Salto abrupto para baixo \\
          \bottomrule
        \end{tabular}
      \end{table}
  \end{itemize}
\end{frame}

\begin{frame}{6. Análise Exploratória (EDA)}
  \begin{itemize}
    \item As séries exibem alta sobreposição de classes, tornando a clusterização desafiadora:
      \begin{itemize}
        \item \emph{Increasing Trend} e \emph{Upward Shift} compartilham dinâmicas de alta no fim da série.
        \item \emph{Decreasing Trend} e \emph{Downward Shift} mostram dinâmicas de queda semelhantes.
        \item \emph{Normal} e \emph{Cyclic} podem se confundir em baixas amplitudes de oscilação.
      \end{itemize}
    \item \textbf{Redução Linear (PCA) vs Não-Linear (t-SNE):}
      \begin{itemize}
        \item O PCA falha em separar as classes Normal e Cyclic de forma isolada.
        \item O t-SNE/UMAP revela 6 nuvens principais, mas sem a ordenação lógica e topológica que o SOM fornece.
      \end{itemize}
  \end{itemize}
\end{frame}

\begin{frame}{7. Pipeline de Trabalho Experimental}
  \begin{itemize}
    \item \textbf{Fluxo de Processamento:}
      \begin{enumerate}
        \item \textbf{Dados:} Leitura e normalização Z-score das séries.
        \item \textbf{Treinamento SOM (Sem Rótulos):} Alimentação do SOM de Kohonen variando o tamanho da grade.
        \item \textbf{Visualização (IntraSOM):} Extração de U-Matrix, Dominância de Classe, Pureza e Component Planes.
        \item \textbf{Avaliação Externa:} Projeção dos rótulos reais sobre os neurônios para medir a acurácia.
        \item \textbf{Métricas:} Cálculo de ARI, NMI, Silhueta, Erro de Quantização, Erro Topográfico.
      \end{enumerate}
  \end{itemize}
\end{frame}

\begin{frame}{8. Tamanhos de Mapa Testados}
  \begin{itemize}
    \item Variamos o tamanho do mapa para analisar o trade-off entre representação geral e granularidade:
      \begin{itemize}
        \item \textbf{5x5 e 7x7:} Grades compactas. Alta densidade de amostras por neurônio, porém causam maior mistura de classes (menor pureza).
        \item \textbf{10x10:} Tamanho padrão equilibrado. Boa separação e baixos erros topológicos.
        \item \textbf{12x12 e 15x15:} Organização intermediária de alta fidelidade.
        \item \textbf{20x20:} Elevada resolução espacial. Risco de neurônios vazios (sem dados projetados), mas com erro de quantização mínimo.
      \end{itemize}
  \end{itemize}
\end{frame}

\begin{frame}{9. Visualização das Fronteiras (U-Matrix)}
  \begin{itemize}
    \item A \textbf{U-Matrix} codifica a distância euclidiana entre o codebook de cada neurônio e seus vizinhos imediatos.
    \item \textbf{Interpretação:}
      \begin{itemize}
        \item Cores escuras/valores baixos indicam alta similaridade (regiões internas dos clusters).
        \item Cores claras/valores altos representam fronteiras de transição (barreiras entre clusters).
      \end{itemize}
    \item Observa-se barreiras claras isolando o grupo \emph{Normal}/\emph{Cyclic} das tendências de queda (\emph{Decreasing}/\emph{Downward}) e de subida (\emph{Increasing}/\emph{Upward}).
  \end{itemize}
\end{frame}

\begin{frame}{10. Mapa de Classes Dominantes}
  \begin{itemize}
    \item Após o treino, mapeamos a classe real majoritária de cada neurônio:
      \begin{itemize}
        \item O SOM organiza os neurônios de forma contínua: classes parecidas ficam geograficamente vizinhas.
        \item A transição de \emph{Increasing Trend} para \emph{Upward Shift} ocorre de forma suave na grade.
        \item As tendências de descida (\emph{Decreasing} e \emph{Downward}) ocupam regiões opostas às tendências de subida.
      \end{itemize}
    \item O mapa exibe uma separação coerente em 6 grandes territórios contíguos na grade.
  \end{itemize}
\end{frame}

\begin{frame}{11. Hits e Pureza dos Neurônios}
  \begin{itemize}
    \item \textbf{Hits (Densidade):} Revela quantas amostras reais caíram em cada neurônio (BMU).
      \begin{itemize}
        \item Neurônios centrais tendem a concentrar mais dados, enquanto neurônios de fronteira têm poucos hits.
      \end{itemize}
    \item \textbf{Pureza:} Proporção de amostras da classe dominante sobre o total de hits do neurônio.
      \begin{itemize}
        \item A maioria dos neurônios ativos possui pureza de 100\% (homogêneos).
        \item Neurônios localizados na divisa física das classes exibem purezas mistas (70\%--85\%), representando as áreas cinzentas do dataset.
      \end{itemize}
  \end{itemize}
\end{frame}

\begin{frame}{12. Component Planes (Instantes de Tempo)}
  \begin{itemize}
    \item O SOM permite projetar pesos sinápticos específicos para atributos individuais (Component Planes).
    \item Analisamos instantes temporais chaves (t=1, t=15, t=30, t=45, t=60):
      \begin{itemize}
        \item Revela a evolução das trajetórias nas regiões da grade.
        \item Permite entender o porquê dos neurônios ativarem: neurônios do canto superior ativam para valores altos de $t$ (tendência crescente), enquanto os inferiores ativam para valores baixos.
      \end{itemize}
  \end{itemize}
\end{frame}

\begin{frame}{13. Métricas Quantitativas do SOM}
  \begin{table}
    \centering
    \begin{tabular}{lcccccc}
      \toprule
      Modelo & ARI & NMI & Silhueta & DB & Erro Quant. & Erro Topo. \\
      \midrule
      SOM 5x5   & 0.6584 & 0.8065 & 0.1548 & 1.7131 & 4.7348 & 0.2400 \\
      SOM 7x7   & 0.6895 & 0.7951 & 0.1876 & 1.5814 & 4.2448 & 0.1289 \\
      SOM 10x10 & 0.6557 & 0.7715 & 0.2042 & 1.5975 & 3.9835 & 0.1500 \\
      SOM 12x12 & 0.6513 & 0.7692 & 0.2089 & 1.6027 & 3.7612 & 0.1483 \\
      SOM 15x15 & 0.6217 & 0.7625 & 0.2040 & 1.6463 & 3.5555 & 0.0659 \\
      SOM 20x20 & 0.5879 & 0.7591 & 0.2629 & 1.2876 & 3.3347 & 0.0459 \\
      \bottomrule
    \end{tabular}
    \caption{Métricas obtidas variando as malhas com IntraSOM.}
  \end{table}
\end{frame}

\begin{frame}{14. Análise de Erros e Pares Difíceis}
  \begin{itemize}
    \item \textbf{Quantização vs Topologia:}
      \begin{itemize}
        \item Conforme o mapa cresce (5x5 para 20x20), o erro de quantização cai de 4.73 para 3.33 (melhor ajuste).
        \item O erro topográfico também cai de 0.24 para 0.04 (melhor preservação da vizinhança).
      \end{itemize}
    \item \textbf{Pares de Confusão Comuns:}
      \begin{itemize}
        \item \emph{Increasing Trend} com \emph{Upward Shift}: Ambas sobem. A diferença é a taxa de crescimento.
        \item \emph{Decreasing Trend} com \emph{Downward Shift}: Ambas descem.
        \item O SOM posiciona esses neurônios na mesma área de transição, validando a coerência científica.
      \end{itemize}
  \end{itemize}
\end{frame}

\begin{frame}{15. Comparação com Baselines de ML}
  \begin{table}
    \centering
    \begin{tabular}{lccccc}
      \toprule
      Método & ARI & NMI & Silhueta & DB & Calinski-Harabasz \\
      \midrule
      \textbf{SOM 10x10} & \textbf{0.6557} & \textbf{0.7715} & 0.2042 & 1.5975 & 208.11 \\
      K-Means Direto     & 0.6120 & 0.7107 & 0.1610 & 1.9582 & 183.15 \\
      PCA + K-Means      & 0.6096 & 0.7109 & 0.1481 & 1.9595 & 181.78 \\
      Clustering Aglom.  & 0.6318 & 0.8008 & 0.2510 & 1.2821 & 209.52 \\
      DBSCAN             & 0.1418 & 0.4019 & 0.1319 & 2.3495 & 55.96 \\
      \bottomrule
    \end{tabular}
  \end{table}
  \begin{itemize}
    \item O SOM supera o K-Means e PCA+K-Means em ARI/NMI e apresenta uma interpretação visual infinitamente superior por ordenar a vizinhança geograficamente.
  \end{itemize}
\end{frame}

\begin{frame}{16. Experimento Textual: Notícias (20 Newsgroups)}
  \begin{itemize}
    \item \textbf{Objetivo:} Mapear notícias textuais na malha hexagonal 10x10 do SOM.
    \item \textbf{Comparação de Vetorização:}
      \begin{enumerate}
        \item \textbf{TF-IDF (Frequencial):} Abordagem Bag-of-Words clássica. ARI = 0.2738.
        \item \textbf{Sentence-BERT (Semântico):} Embeddings contextuais densos. ARI = 0.5873.
      \end{enumerate}
    \item \textbf{Análise:}
      \begin{itemize}
        \item O TF-IDF falha quando há sinônimos ou palavras diferentes para o mesmo tema.
        \item O Sentence-BERT mapeia significado e conceitos abstratos. O SOM treinado sobre SBERT agrupa com clareza as quatro temáticas e permite inferência em tempo real confiável.
      \end{itemize}
  \end{itemize}
\end{frame}

\begin{frame}{17. Aplicações em Gestão Pública, IA e Direito}
  \begin{itemize}
    \item \textbf{Triagem e Agrupamento de Processos Judiciais:} Agrupar petições iniciais semelhantes para análise em lote (Visual Law e Jurimetria).
    \item \textbf{Jurisprudência Inteligente:} Criar mapas visuais interativos de ementas e decisões para identificar tendências de julgamento de tribunais.
    \item \textbf{Monitoramento Temático de Notícias:} Acompanhamento em tempo real de tópicos sensíveis ou tendências de imprensa na segurança e inteligência.
    \item \textbf{Gestão Documental (Memória Técnica):} Organização não-supervisionada de relatórios, memorandos e memórias técnicas institucionais da administração pública.
  \end{itemize}
\end{frame}

\begin{frame}{18. Conclusão}
  \begin{itemize}
    \item \textbf{Conclusão:} O SOM de Kohonen se provou uma ferramenta robusta e cientificamente rica tanto para análise temporal quanto textual.
    \item \textbf{Resultados Alcançados:}
      \begin{itemize}
        \item Preservação topográfica excelente em séries temporais (erros topográficos mínimos).
        \item Superioridade gritante das representações semânticas (SBERT) para clusterização de textos.
        \item Desenvolvimento de um demonstrador funcional interativo (Dashboard Tokyo Night) que atende a todos os critérios e métricas exigidos.
      \end{itemize}
    \item \textbf{Trabalhos Futuros:} Incorporação de mapas toroidais e visualização de trajetórias temporais animadas.
  \end{itemize}
\end{frame}

\end{document}
